% TEKST %
% Setter T1-fontkoding slik at spesialtegn (æ, ø, å) håndteres riktig ved bokstavinndeling og søk/kopiering
\usepackage[T1]{fontenc} 

% PAKKER %
% Her importerer jeg pakker, om en ikke har bruk for å importere med spesielle instillinger så kan de bare legges i denne listen
\usepackage{quoting, tabularx, fancyhdr, ragged2e, fullpage, float, booktabs, subcaption}
% xltabular gir longtable med X-kolonner (tabeller som kan brytes over flere sider)
\usepackage{xltabular}
% L er en venstrejustert X-kolonne, som unngår store mellomrom i smale tabellceller
\newcolumntype{L}{>{\raggedright\arraybackslash}X}
% Tabelltekster plasseres over tabellen
\captionsetup[table]{position=top}
% placeins med section-flagg tvinger floats til å bli i sin egen seksjon
\usepackage[section]{placeins}

%% XCOLOR %%
% xcolor behøver table innstillingen for å kunne fargere celler i tabeller, og xcdraw for enkelte funksjoner i TablesGenerator.com
\usepackage[table,xcdraw]{xcolor}

%% TODONOTES %%
% Tar i bruk pakker her som har behov for noen andre innstillinger enn hva som er satt fra før.
% todonotes lar en legge inn kommentarer/oppgaver i margen, og colorinlistoftodos gir fargekoder i oversikten over todos
\usepackage[colorinlistoftodos]{todonotes}

%% SUBFILES %%
% subfiles lar en dele dokumentet opp i flere filer som kan kompileres hver for seg
\usepackage{subfiles}

%% TOCBIBIND %%
% Sørger for at innholdsfortegnelse, figurliste og tabelliste tas med i selve innholdsfortegnelsen (her: med unntak av disse via not-flaggene)
\usepackage[nottoc,notlof,notlot]{tocbibind}

%% KILDEHÅNDTERING %%
% Her så tar en i bruk biblatex pakken, som håndterer kildene som blir importert av zotero
\usepackage[backend=biber,style=ieee,sorting=none]{biblatex}

% Her legger en til biblatex sin bibliografifil
\addbibresource{bib/references.bib}

% Lar lange URL-er brytes også etter små og store bokstaver, så de ikke går ut i margen
\setcounter{biburllcpenalty}{7000}
\setcounter{biburlucpenalty}{8000}
% Ingen strekk inne i URL-er (ellers blir det "https : / / www .")
\biburlbigskip=0mu \biburlnumskip=0mu \biburlucskip=0mu \biburllcskip=0mu
% Litt ekstra strekk i kildelisten, så ingen linjer går ut i margen
\AtBeginBibliography{\setlength{\emergencystretch}{2em}}

%% MINTED %%
% Dette er i bruk for å kunne få med kode i rapporten
\usepackage{minted}

% Definerer en lysegrå farge som kan brukes som bakgrunn på kodeblokker
\definecolor{LightGray}{gray}{0.9}

%% CSQUOTES %%
% csquotes håndterer anførselstegn riktig, og kreves av biblatex for korrekt sitering
\usepackage{csquotes}

%% GRAPHICX %%
% Denne er en del av graphicx pakken, denne forteller overleaf at den skal se etter figurer og andre bilder i mappen images.
\usepackage{graphicx}
% pdflscape gir liggende sider (roteres også i PDF-leseren), afterpage lar
% en slik side komme etter den nåværende siden uten å avbryte teksten
\usepackage{pdflscape, afterpage}
% \landscapefigure{fil}{tekst}{label}: diagram på en egen liggende side etter
% den nåværende siden, nesten hele arkets bredde (marger på 1 cm)
\newcommand{\landscapefigure}[3]{\afterpage{%
  \begin{landscape}
  \begin{figure}[p]
    \centering
    \makebox[\linewidth][c]{\includegraphics[width=\dimexpr\paperheight-2cm\relax,height=0.9\textheight,keepaspectratio]{#1}}
    \caption{#2}
    \label{#3}
  \end{figure}
  \end{landscape}}}
% \landscapefigurenopage{fil}{tekst}{label}: samme, men uten sidetall (for
% diagrammer som går over bunnteksten)
\newcommand{\landscapefigurenopage}[3]{\afterpage{%
  \begin{landscape}
  \thispagestyle{empty}
  \begin{figure}[p]
    \centering
    \makebox[\linewidth][c]{\includegraphics[width=\dimexpr\paperheight-2cm\relax,height=0.9\textheight,keepaspectratio]{#1}}
    \caption{#2}
    \label{#3}
  \end{figure}
  \end{landscape}}}
% \wideimage{fil}: bilde som går litt ut i margene (arkets bredde minus 2 cm)
\newcommand{\wideimage}[1]{\makebox[\textwidth][c]{%
  \includegraphics[width=\dimexpr\paperwidth-2cm\relax,height=0.8\textheight,keepaspectratio]{#1}}}
\graphicspath{{images/}{../images/}} 

%% HYPERREF %%
% Vi vil ikke at lienker skal være klikkbare og må derfor se til at den har hidelinks, 
\usepackage[hidelinks]{hyperref}

%% MARGER %%
% Denne setter hvor stor margene skal være i dokumentet, dette er for å kunne gjøre kommentarer med f.eks \todo
\setlength{\marginparwidth}{2cm}